\documentclass[11pt]{scrartcl}

\usepackage{amsmath,amsthm,amssymb,mathtools}
\usepackage{hyperref}
\usepackage{lastpage}
\usepackage{multicol}

\theoremstyle{plain} 
\newtheorem{theorem}{Theorem}[section]
\newtheorem{lemma}[theorem]{Lemma}
\newtheorem*{question}{Question}
\theoremstyle{definition}
\newtheorem*{definition}{Definition} 
\theoremstyle{remark}
\newtheorem*{remark}{Remark}

\usepackage[none]{hyphenat}
\setlength{\parindent}{0pt}
% setlength{\emergencystretch}{3em}

\renewcommand\labelitemi{$\vcenter{\hbox{\tiny$\bullet$}}$} 

\usepackage[margin=1in]{geometry}
\usepackage{fancyhdr}
\pagestyle{fancy}
\lhead{ESC194: Calculus I} 
\rhead{\thepage}
\cfoot{}

\title{ESC194: Calculus I \\[0.1in] Exam Review}
\author{Brendan Tang}
\date{Fall 2026}

\begin{document}

\maketitle
\thispagestyle{empty}

\tableofcontents

\newpage
\section{Limits and Continuity}
\subsection{The Precise Definition of a Limit}

The precise definition of a limit is as follows.

\begin{definition}
    If $\forall\epsilon>0$, $\exists\delta>0$ such that $\lvert x-a \rvert<\delta\implies
    \lvert f(x)-L \rvert<\epsilon$, then \[ \lim_{x \to a}f(x)=L. \]
\end{definition}

For one-sided limits, $x$ is only approaching $a$ from one side. Thus, we define the left-handed
limit as follows.

\begin{definition}
    If $\forall\epsilon>0$, $\exists\delta>0$ such that $a-\delta<x<a\implies \lvert f(x)-L \rvert
    <\epsilon$, then \[ \lim_{ x \to a^- }f(x)=L. \]
\end{definition}

As a number becomes arbitrarily large, it eventually surpasses every real number. Thus, we define
limits to or approaching infinity as follows.

\begin{definition}
    If $\forall M>0$, $\exists\delta>0$ such that $0<\lvert x-a \rvert<\delta\implies f(x)>M$, then
    \[ \lim_{ x \to a }f(x) \to \infty. \]
    Conversely, if $\forall\epsilon>0$, $\exists M>0$ such that $x>M\implies \lvert f(x)-L \rvert
    <\epsilon$, then \[ \lim_{ x \to \infty }f(x)=L. \]
\end{definition}

\begin{question}
    How can we rigorously define right-handed and negative infinite limits? What about limits 
    involving both concepts?
\end{question}

Here are the general steps to a rigorous proof.

\begin{enumerate}
    \item Given $\epsilon>0$, find $\delta>0$. Start with $\lvert f(x)-L \rvert<\epsilon$ and manipulate
    it into the form $\lvert x-a \rvert<\delta_{0}$. This step may require bounding $\delta\leq n$.
    \item \begin{proof} 
        Given $\epsilon$, let $\delta=min(n,\delta_{0})$. Then, show that $\lvert x-a 
        \rvert<\delta\implies \lvert f(x)-L \rvert<\epsilon$, as desired. 
    \end{proof}
\end{enumerate}

\subsection{Limit Theorems}

Here are some important limit theorems.

\begin{theorem}
    If $\lim_{ x \to a }f(x)$ and $\lim_{ x \to a }g(x)$ exist, then
    \begin{itemize}
        \item (Constant LT) $\lim_{ x \to a }c=c$
        \item (Sum LT) $\lim_{ x \to a }[f(x)+g(x)]=\lim_{ x \to a }f(x)+\lim_{ x \to a }g(x)$
        \item (Product LT) $\lim_{ x \to a }[f(x)\cdot g(x)]=\lim_{ x \to a }f(x)\cdot 
        \lim_{ x \to a }g(x)$
        \item (Constant multiple LT) $\lim_{ x \to a }cf(x)=c\lim_{ x \to a }f(x)$
        \item (Power LT) $\lim_{ x \to a }[f(x)]^n=[\lim_{ x \to a }f(x)]^n$
        \item (Root LT) $\lim_{ x \to a }\sqrt[n]{ f(x) }=\sqrt[n]{ \lim_{ x \to a }f(x) }$
    \end{itemize}
\end{theorem}

\begin{theorem}[Direct Substitution]
    If $f$ is continuous at $a$, then \[ \lim_{ x \to a } f(x)=a. \]
\end{theorem}

The following theorem is often used to evaluate trigonometric limits.

\begin{theorem}[Squeeze Theorem]
    If $g(x)<f(x)<h(x)$ and $\lim_{ x \to a }g(x)=\lim_{ x \to a }h(x)=L$, then 
    \[ \lim_{ x \to a } f(x)=L. \]
\end{theorem}

When solving difficult trigonometric limits, it may help to find $\frac{{\sin x}}{x}$ terms.

\subsection{Continuity Theorems}

Here is the definition for continuity.

\begin{definition}
    $f$ is continuous at a point $a$ if $\lim_{ x \to a }f(x)=f(a)$. 
\end{definition}

\begin{question} 
    How can we define continuity from the side? What about continuity on an open or closed interval?
\end{question}

Here are some continuity theorems.

\begin{theorem}
    If $f$ and $g$ are continuous at $a$, then \[ f+g,\ fg,\ cf \] are also continuous at $a$.
\end{theorem}

\begin{theorem}
    The following types of functions are continuous on their domains 
    \begin{multicols}{3}
        \begin{itemize}
            \item polynomial
            \item exponential
            \item rational
            \item logarithmic
            \item root
            \item trigonometric
        \end{itemize}
    \end{multicols}
\end{theorem}

\begin{remark}
    It turns out that all elementary functions are continuous on their domains (why?).
\end{remark}

\newpage
\section{Derivatives}
\subsection{Derivative Theorems}

We begin with some important definitions.

\begin{definition}
    The derivative of $f$ is defined as \[ f'(x)\coloneqq \lim_{ h \to 0 } \frac{{f(x+h)-f(x)}}{h}. \]
    $f$ is differentiable at a point $a$ if $f'(a)$ exists.
\end{definition}

\begin{question} 
    How can we define differentiability from the side? What about differentiability on an open or
    closed interval?
\end{question}

From the definition of the derivative, we see that differentiability at $a$ requires $\lim_{ h \to 0 } 
[f(a+h)-f(a)]=0$ (why?). This implies that $f$ is continuous at $a$. Thus, we arrive at the following theorem.

\begin{theorem}
    If $f$ is differentiable at $a$, then $f$ is continuous at $a$.
\end{theorem}

Be careful, as this theorem does not hold for the converse. That is, continuity at $a$ does not imply
differentiability at $a$. \\[-0.5em]

Here are some useful differentiation theorems.

\begin{theorem}
    If $f$ and $g$ are differentiable, then
    \begin{itemize}
        \item (Constant DT) $f(x)=c\implies f'(x)=0$
        \item (Sum DT) $(f+g)'=f'+g'$
        \item (Product DT) $(fg)'=f'g+fg'$
        \item (Reciprocal DT) $\left( \frac{1}{f} \right)'=-\frac{f'}{f^{2}}$
        \item (Quotient DT) $\left( \frac{f}{g} \right)'= \frac{{f'g-fg'}}{g^{2}}$
        \item (Power DT) $f(x)=x^n\implies f'(x)=nx^{n-1}$
        \item (Chain rule) $(f\circ g)'=(f'\circ g)\cdot g'$
    \end{itemize}
\end{theorem}

\subsubsection{Existence Theorems}

These theorems are very important, serving as stepping stones for many proofs that will be encountered in
this course. While their proofs will not be provided here, it is highly encouraged to find and read through 
them yourself.

\begin{theorem}[Intermediate Value Theorem]
    If $f\colon[a,b]\to \mathbb{R}$ is continuous, then for every $N\in\bigl(f(a),f(b)\bigr)$, there exists 
    $c\in(a,b)$ such that $f(c)=N$.
\end{theorem}

\begin{theorem}[Extreme Value Theorem]
    If $f\colon[a,b]\to \mathbb{R}$ is continuous, then there exists $c,d\in[a,b]$ such that $f$ attains an 
    absolute maximum value $f(c)$ and an absolute minimum value $f(d)$.
\end{theorem}

\begin{theorem}[Mean Value Theorem] \label{thm:mvt}
    If $f\colon[a,b]\to \mathbb{R}$ is continuous and differentiable on $(a,b)$, there exists $c\in(a,b)$ 
    such that \[ f'(c)= \frac{{f(b)-f(a)}}{b-a}. \]
\end{theorem}

By using Theorem \ref{thm:mvt} with $f(a)=f(b)$, we get a special case.

\begin{theorem}[Rolle's Theorem] \label{thm:rolle}
    If $f\colon[a,b]\to \mathbb{R}$ is continuous, differentiable on $(a,b)$, and $f(a)=f(b)$, there exists 
    $c\in(a,b)$ such that $f'(c)=0$.
\end{theorem}

\begin{remark}
    Although Theorem \ref{thm:rolle} is presented here as a special case of Theorem \ref{thm:mvt}, 
    chronologically the reverse is true. In fact, Theorem \ref{thm:rolle} is the core tool used to prove 
    Theorem \ref{thm:mvt}, which serves as its generalization.
\end{remark}

\subsubsection{Interlude (Approximation)}

On a bit of an unrelated note, we can use differentials to approximate a function $f$ by the formula
\[ f(x+\Delta x)\approx f(x)+f'(x)\Delta x. \]
This has been included for the sake of completeness, as it has been tested before.

\subsection{Curve Sketching}

By Fermat's theorem on stationary points, we arrive at the following definition.

\begin{definition}
    $\bigl(c,f(c)\bigr)$ is a critical point on $f$ if $f'(c)=0$ or $f'(c)$ does not exist.
\end{definition}

Here are the general steps to curve sketching.

\begin{enumerate}
    \item From $f$:
    \begin{itemize}
        \item Find the domain and range 
        \item Find the $x$ and $y$ intercepts
        \item Find any endpoints or removable discontinuities, if they exist
        \item Find any asymptotes and their behaviour, if they exist
        \item Determine if $f$ has any symmetry
        \item Determine the end behaviour of $f$
    \end{itemize}
    \item From $f'$: 
    \begin{itemize}
        \item Find all critical points on $f$
        \item Determine the intervals of increase/decrease
        \item Find vertical tangents or cusps, if they exist
    \end{itemize}
    \item From $f''$: 
    \begin{itemize}
        \item Find local maxima/minima and any inflection points, if they exist
        \item Determine the intervals of concavity
    \end{itemize}
\end{enumerate}

\subsection{Applications of Derivatives}

One application of derivatives is related rates. Related rates are problems where two or more quantities change,
but are still linked by an equation. To solve a related rates problem, find an equation that directly relates 
the variables and apply the chain rule. \\[-0.5em]

Another application of derivatives is optimization. To solve an optimization problem, reduce the problem to an
equation in one variable, then evaluate at endpoints and critical points.

\newpage
\section{Integrals}
\subsection{Summation}

Here are some important summation theorems.

\begin{theorem}
    If $\sum_{k=1}^n a_k$ and $\sum_{k=1}^n b_k$ exist, then
    \begin{itemize}
        \item $\sum_{k=1}^n [a_k+b_k]=\sum_{k=1}^n a_k+\sum_{k=1}^n b_k$
        \item $\sum_{k=1}^n [a_k\cdot b_k]=\sum_{k=1}^n a_k\cdot\sum_{k=1}^n b_k$
        \item $\sum_{k=1}^n ca_k=c\sum_{k=1}^n a_k$
        \item $\sum_{k=1}^n c=nc$
        \item $\sum_{k=1}^n k=\frac{n(n+1)}{2}$
        \item $\sum_{k=1}^n k^{2}=\frac{n(n+1)(2n+1)}{6}$
        \item $\sum_{k=1}^n k^{3}=\left[ \frac{n(n+1)}{2} \right]^{2}$
    \end{itemize}
\end{theorem}

There is also a general formula for the sum of the first $n$ $p$th powers.

\begin{theorem}[Faulhaber's Formula]
    Let $p$ be a positive integer. Then
    \[ \sum_{k=1}^n k^p=\frac{1}{p+1}\sum_{r=0}^p B_{r}^+ n^{p+1-r},\]
    where $B_{r}^+$ are the second Bernoulli numbers. 
\end{theorem}

\subsubsection{The Riemann Integral}

We begin with a very important definition.

\begin{definition}
    Let $f\colon[a,b]\to \mathbb{R}$ be continuous and $P=\{x_{0},x_{1},\dots,x_{n}\}$ be a partition of $[a,b]$ where
    \[ a=x_{0}<x_{1}<\dots<x_{n}=b. \]
    Then the Riemann sum $S$ of $f$ over $[a,b]$ with a partition $P$ is defined as
    \[ S=\sum_{i=1}^n f(x_{i}^*)\Delta x_{i}, \]
    where $\Delta x_{i}=x_{i}-x_{i-1}$ and $x_{i}^*\in[x_{i-1},x_{i}]$. 
\end{definition}

From this, we derive the definition of the definite integral.

\begin{definition}
    Let $f\colon[a,b]\to \mathbb{R}$ be continuous and $P=\{x_{0},x_{1},\dots,x_{n}\}$ be a partition of $[a,b]$ where
    \[ a=x_{0}<x_{1}<\dots<x_{n}=b. \]
    Then the definite integral of $f$ over $[a,b]$ with a norm $\lVert P \rVert=\max\{\Delta x_{i}\}$ is defined as
    \[ \int_{a}^{b} f(x) \, dx \coloneqq \lim_{ \lVert P \rVert  \to 0 } \sum_{i=1}^n f(x_{i}^*)\Delta x_{i} \]
    if the limit exists, where $\Delta x_{i}=x_{i}-x_{i-1}$ and $x_{i}^*\in[x_{i-1},x_{i}]$. If the limit does exist, then $f$ is integrable on $[a,b]$.
\end{definition}

From this definition of the integral, we see that the difference between the upper and lower Riemann sums needs to be arbitrarily small. If $f$ is continuous on $[a,b]$, then we can make this difference arbitrarily small by choosing a sufficiently fine partition. Thus, we arrive at the following theorem.

\begin{theorem}
    If $f$ is continuous (or piecewise continuous) on $[a,b]$, then $f$ is Riemann integrable on $[a,b]$. 
\end{theorem}

When evaluating a Riemann sum, we often select regular partitions and take the right Riemann sum. Thus, we get $\Delta x_{i}=\frac{{b-a}}{n}$, $x_{i}^*=a+i \frac{{b-a}}{n}$, and so the integral can be written as
\[ \int_{a}^{b} f(x) \, dx =\lim_{ n \to \infty } \sum_{i=1}^n f\left( a+i \frac{{b-a}}{n} \right) \frac{{b-a}}{n}. \]

\subsection{Integral Theorems}

Here are some theorems for definite integrals.

\begin{theorem}
    Let $f,g\colon[a,b]\to \mathbb{R}$ be integrable. Then
    \begin{itemize}
        \item $\int_{a}^{b} [f(x)+g(x)] \, dx = \int_{a}^{b} f(x) \, dx+\int_{a}^{b} g(x) \, dx$
        \item $\int_{a}^{b} cf(x) \, dx=c\int_{a}^{b} f(x) \, dx$
        \item $\int_{a}^{b} c \, dx=c(b-a)$
        \item $\int_{a}^{b} f(x) \, dx=\int_{a}^{c} f(x) \, dx+\int_{c}^{b} f(x) \, dx$
    \end{itemize} 
\end{theorem}

We also have the order properties of integrals. 

\begin{theorem}
    Let $f,g\colon[a,b]\to \mathbb{R}$ be integrable. Then
    \begin{itemize}
        \item For all $f$, 
        \[ \int_{a}^{b} f(x) \, dx =-\int_{b}^{a} f(x) \, dx  \]
        \item If $f(x)\leq g(x)$ for all $x \in[a,b]$, then 
        \[ \int_{a}^{b} f(x) \, dx \leq \int_{a}^{b} g(x) \, dx  \]
        \item If $m\leq f(x)\leq M$ for all $x \in[a,b]$, then
        \[ m(b-a)\leq \int_{a}^{b} f(x) \, dx \leq M(b-a) \]
        \item For all $f$,
        \[ \left\lvert  \int_{a}^{b} f(x) \, dx   \right\rvert \leq \int_{a}^{b} \lvert f(x) \rvert  \, dx  \]
    \end{itemize}
\end{theorem}

Here is the first fundamental theorem of calculus.

\begin{theorem}[First Fundamental Theorem of Calculus]
    For a continuous function $f$, an antiderivative $F$ can be obtained as the integral of $f$ over an interval with a variable upper bound. That is,
    \[ F(x)=\int_{a}^{x} f(t) \, dt  \]
    and
    \[ F'(x)=\frac{d}{dx}\left[ \int_{a}^{x} f(t) \, dt  \right]=f(x). \]
\end{theorem}

If the upper bound is a function $g(x)$, we can use the chain rule to get
\[ F'(x)=\frac{d}{dx}\left[ \int_{a}^{g(x)} f(t) \, dt  \right] = f(g(x))\cdot g'(x). \]

Here is the second fundamental theorem of calculus.

\begin{theorem}[Second Fundamental Theorem of Calculus]
    The integral of $f$ over a fixed interval is equal to the change of any antiderivative $F$ between the ends of the interval. That is,
    \[ \int_{a}^{b} f(x) \, dx =F(b)-F(a). \]
\end{theorem}

\subsection{Integration Techniques}

We can use the results of common derivatives to find integrals. Here are some standard integrals:

\begin{itemize}
    \begin{multicols}{2}
        \item $\int u^n \, du = \frac{u^{n+1}}{n+1}+C$, $n\neq 1$
        \item $\int \frac{du}{u} =\ln \lvert u \rvert+C$
        \item $\int e^u \, du =e^u+C$
        \item $\int a^u \, du =\frac{a^u}{\ln a}+C$
        \item $\int \sin u \, du =-\cos u+C$
        \item $\int \cos u \, du =\sin u+C$
        \item $\int \sec ^{2}u \, du = \tan u+C$
        \item $\int \csc ^{2}u \, du = -\cot u+C$
        \item $\int \sec u\tan u \, du = \sec u+C$
        \item $\int \csc u\cot u \, du = -\csc u+C$
        \item $\int \tan u \, du = \ln \lvert \sec u \rvert+C$
        \item $\int \cot u \, du = \ln \lvert \sin u \rvert+C$
        \item $\int \sec u \, du = \ln \lvert \sec u+\tan u \rvert+C$
        \item $\int \csc u \, du = \ln \lvert \csc u-\cot u \rvert+C$
        \item $\int \frac{du}{\sqrt{ a^{2}-u^{2} }}=\sin^{-1} \frac{u}{a}+C$, $a>0$
        \item $\int \frac{du}{a^{2}+u^{2}}=\frac{1}{a} \tan^{-1} \frac{u}{a}+C$
        \item $\int \frac{du}{u\sqrt{ u^{2}-a^{2} }}=\frac{1}{a} \sec^{-1} \frac{u}{a}+C$
        \item $\int \frac{du}{u\sqrt{ a^{2}-u^{2} }}=\frac{1}{2a} \ln \left\lvert  \frac{{u+a}}{u-a}  \right\rvert+C$
        \item $\int \frac{du}{u\sqrt{ u^{2}-a^{2} }}=\frac{1}{2a} \ln \left\lvert  \frac{{u-a}}{u+a}  \right\rvert+C$
    \end{multicols}
\end{itemize}

Sometimes an integral will not be in a form which we can easily integrate. Another technique we can apply is substitution, which will reverse the chain rule. Let $u=g(x)$, then
\[ \int f(g(x))g'(x) \, dx = \int f(u) \, du. \]

Some integrals with square roots can be simplified by taking advantage of the pythagorean identity. To do this, we can make a trigonometric substitution:

\begin{itemize}
    \item $\sqrt{ a^{2}-b^{2}x^{2} }\implies$choose $x=\frac{a}{b}\sin \theta$
    \item $\sqrt{ a^{2}+b^{2}x^{2} }\implies$choose $x=\frac{a}{b}\tan \theta$
    \item $\sqrt{ b^{2}x^{2}-a^{2} }\implies$choose $x=\frac{a}{b}\sec \theta$
\end{itemize}

\begin{question}
    What will each of these expressions simplify into?
\end{question}

Sometimes we may have a nonzero $x$ term inside the square root. We can turn it into one of the above forms by completing the square. \newpage

Here are some more theorems that may be helpful for integration.

\begin{theorem}
    If $f$ is an odd function and integrable on $[-a,a]$, then
    \[ \int_{-a}^{a} f(x) \, dx =0. \]
    If $f$ is an even function and integrable on $[-a,a]$, then
    \[ \int_{-a}^{a} f(x) \, dx =2\int_{0}^{a} f(x) \, dx . \]
\end{theorem}

\begin{theorem}[King's Rule]
    Let $f\colon[a,b]\to \mathbb{R}$ be integrable. Then
    \[ \int_{a}^{b} f(x) \, dx =\int_{a}^{b} f(a+b-x) \, dx . \]
\end{theorem}

A rational trigonometric integral can be transformed into an ordinary rational integral by the Weierstrass substitution, $t=\tan \frac{x}{2}$. This substitution implies the following:

\begin{itemize}
    \begin{multicols}{2}
        \item $\sin x=\frac{2t}{1+t^{2}}$
        \item $\cos x=\frac{{1-t^{2}}}{1+t^{2}}$
        \item $\tan x=\frac{2t}{1-t^{2}}$
        \item $dx=\frac{2}{1+t^{2}} \, dt$
    \end{multicols}
\end{itemize}

\subsection{Applications of Integrals}

One application of integrals is finding the area between two curves. This is done similarly to taking the area under a curve and the $x$-axis, except now we take the definite integral
\[ \int_{a}^{b} \lvert f(x)-g(x) \rvert  \, dx . \]

\end{document}